\chapter{Catálogo de errores}
\label{ap:errores}

Esta es la lista de las cosas que salen mal, con el síntoma por el que se
manifiestan. Está ordenada por el momento en que aparecen.

\section{En el estado estacionario}

\paragraph{La distribución no suma uno} Casi siempre es la lotería: masa que se
escapa por el extremo de la grilla porque la política excede el último punto y no
se trunca. Síntoma: la masa total deriva lentamente hacia abajo con las
iteraciones.

\paragraph{Hay masa apreciable en el último punto de la grilla} La grilla está
atando. En un modelo de hogares suele bastar con extenderla. En un modelo con
acumulación y retornos lineales puede indicar algo peor: que la condición de
crecimiento $\beta\varsigma R<1$ está violada y la distribución verdaderamente
diverge.

\paragraph{El presupuesto agregado no cierra} Error de contabilidad en el paso
hacia atrás. Hay que verificar punto por punto que
$c + a' = (1+r)a_- + we$ sobre toda la grilla, no solo en el agregado.

\paragraph{La iteración hacia atrás no converge} Revisar que la condición de
transversalidad implícita se cumpla ($\beta R<1$ en el caso estándar), y que el
valor inicial de la variable de continuación sea finito y positivo en todos los
puntos.

\section{En los jacobianos}

\paragraph{No coinciden con el método directo} Las causas, por frecuencia:
\begin{enumerate}[leftmargin=1.4em,itemsep=1pt]
\item orden invertido de los operadores en el paso de expectativa (ver apéndice
\ref{ap:codigo});
\item la variable de continuación se reinicia dentro del bucle sobre $s$ en el
paso 1, en lugar de arrastrarse;
\item desfase de un índice en la recursión del paso 4;
\item el estado estacionario no convergió lo suficiente;
\item hay un rezago del insumo dentro del bloque heterogéneo (ver capítulo
\ref{cap:extensiones}).
\end{enumerate}

\paragraph{La primera fila no es cero para una variable de acervo} Error de
alineación temporal: en algún lado se está usando la distribución de $t+1$ donde
corresponde la de $t$.

\paragraph{El jacobiano empeora al reducir $h$} El estado estacionario está mal
convergido y su ruido domina la señal de la derivada.

\paragraph{El jacobiano tiene entradas enormes cerca de la diagonal} Puede ser
real --- restricciones de liquidez fuertes --- o puede ser que la política tenga
un salto sobre la grilla. Vale la pena graficar la política antes y después de
perturbar.

\section{En el equilibrio general}

\paragraph{\texttt{LinAlgError: singular matrix}} Casi siempre un objetivo
redundante: se impusieron todas las condiciones de vaciado y la ley de Walras
hace a una de ellas combinación lineal de las demás. También puede ser un objetivo
insensible a toda incógnita.

\paragraph{La solución es enorme o cambia de signo al cambiar $T$} El sistema
está casi singular. Mirar el número de condición de $\mathbf{H}_U$.

\paragraph{La respuesta no decae hacia $T$} El horizonte es corto. La solución
truncada impone implícitamente que todo vuelva al estado estacionario en $T$, y
si no lo hace, el error se propaga hacia atrás.

\paragraph{Una variable predeterminada responde en impacto} Error en el grafo:
probablemente un rezago mal puesto en un bloque simple.

\section{En el modelo}

\paragraph{El estado estacionario está sobre un quiebre} Si una función $\min$ o
$\max$ se evalúa exactamente en su punto de quiebre, la derivada no existe. Hay
que recalibrar. Si está \emph{cerca} del quiebre, el modelo es formalmente
diferenciable pero la linealización no es informativa, y conviene decirlo.

\paragraph{La distribución no afecta a nada} Si todos los productos del bloque son
lineales en el estado endógeno, el bloque colapsa a un conjunto de agentes
representativos. Síntoma: los jacobianos no cambian al modificar la forma de la
distribución manteniendo su media. Ver capítulo \ref{cap:macrofinanzas}.

\paragraph{Los resultados cambian mucho con el número de puntos de la grilla}
Error de discretización, no del método. Suele resolverse concentrando puntos
donde la política tiene curvatura, no simplemente agregando más.
